% ======================================================================
%  Preámbulo común de los reportes de Teoría de la Computación
%  Compilar con XeLaTeX + Biber:  latexmk -xelatex main.tex
% ======================================================================
\usepackage{fontspec}
\IfFontExistsTF{Times New Roman}
  {\setmainfont{Times New Roman}}
  {\setmainfont{TeX Gyre Termes}}
\setmonofont{DejaVu Sans Mono}[Scale=0.82]
\newfontfamily\simbolos{DejaVu Sans}[Scale=0.92]

\usepackage[spanish,mexico,es-noshorthands]{babel}
\usepackage{csquotes}

\usepackage[letterpaper,margin=2.54cm,headheight=15pt]{geometry}
\usepackage{setspace}
\doublespacing

\usepackage{amsmath,amssymb}
\usepackage{graphicx}
\usepackage[table]{xcolor}
\usepackage{booktabs,longtable,array,calc,multirow}
\usepackage{caption}
\captionsetup{font={small,singlespacing},labelfont=bf,labelsep=period}
\usepackage{enumitem}
\setlist{nosep,leftmargin=*}
\usepackage{tikz}
\usepackage{fancyhdr}
\usepackage{titlesec}
\usepackage[style=apa,backend=biber]{biblatex}
\usepackage[hidelinks]{hyperref}

% --- Símbolos Unicode que la fuente de texto no incluye ------------------
\usepackage{newunicodechar}
\newunicodechar{⁰}{\textsuperscript{0}}
\newunicodechar{⁴}{\textsuperscript{4}}
\newunicodechar{⁵}{\textsuperscript{5}}
\newunicodechar{⁶}{\textsuperscript{6}}
\newunicodechar{⁷}{\textsuperscript{7}}
\newunicodechar{⁸}{\textsuperscript{8}}
\newunicodechar{⁹}{\textsuperscript{9}}
\newunicodechar{⁺}{\textsuperscript{+}}
\newunicodechar{ⁿ}{\textsuperscript{\textit{n}}}
\newunicodechar{ᵏ}{\textsuperscript{\textit{k}}}
\newunicodechar{ᴿ}{\textsuperscript{\textit{R}}}
\newunicodechar{₀}{\textsubscript{0}}
\newunicodechar{₁}{\textsubscript{1}}
\newunicodechar{₂}{\textsubscript{2}}
\newunicodechar{₃}{\textsubscript{3}}
\newunicodechar{∅}{\ensuremath{\varnothing}}
\newunicodechar{∈}{\ensuremath{\in}}
\newunicodechar{∉}{\ensuremath{\notin}}
\newunicodechar{∣}{\ensuremath{\mid}}
\newunicodechar{∩}{\ensuremath{\cap}}
\newunicodechar{∪}{\ensuremath{\cup}}
\newunicodechar{⊆}{\ensuremath{\subseteq}}
\newunicodechar{⋃}{\ensuremath{\bigcup}}
\newunicodechar{⋯}{\ensuremath{\cdots}}
\newunicodechar{≤}{\ensuremath{\le}}
\newunicodechar{≥}{\ensuremath{\ge}}
\newunicodechar{≠}{\ensuremath{\neq}}
\newunicodechar{→}{\ensuremath{\to}}
\newunicodechar{↦}{\ensuremath{\mapsto}}
\newunicodechar{×}{\ensuremath{\times}}
\newunicodechar{·}{\ensuremath{\cdot}}
\newunicodechar{𝒫}{\ensuremath{\mathcal{P}}}
\newunicodechar{✓}{{\simbolos ✓}}
\newunicodechar{✗}{{\simbolos ✗}}
\newunicodechar{↔}{\ensuremath{\leftrightarrow}}
\newunicodechar{⇔}{\ensuremath{\Leftrightarrow}}
\newunicodechar{∋}{\ensuremath{\ni}}
\newunicodechar{⊥}{\ensuremath{\bot}}

\newunicodechar{ᵐ}{\textsuperscript{\textit{m}}}
\newunicodechar{ᵖ}{\textsuperscript{\textit{p}}}
\newunicodechar{ᵢ}{\textsubscript{\textit{i}}}
\newunicodechar{ⱼ}{\textsubscript{\textit{j}}}
\newunicodechar{ₓ}{\textsubscript{\textit{x}}}
\newunicodechar{⁻}{\textsuperscript{−}}
\newunicodechar{≡}{\ensuremath{\equiv}}
\newunicodechar{¹}{\textsuperscript{1}}
\newunicodechar{ᵃ}{\textsuperscript{\textit{a}}}
\newunicodechar{ᵇ}{\textsuperscript{\textit{b}}}
\newunicodechar{ᶜ}{\textsuperscript{\textit{c}}}
\newunicodechar{∃}{\ensuremath{\exists}}
\newunicodechar{∀}{\ensuremath{\forall}}

\newunicodechar{ˣ}{\textsuperscript{\textit{x}}}
\newunicodechar{ⁱ}{\textsuperscript{\textit{i}}}
\newunicodechar{ₙ}{\textsubscript{\textit{n}}}
\newunicodechar{⟶}{\ensuremath{\longrightarrow}}
\newunicodechar{⇒}{\ensuremath{\Rightarrow}}
\newunicodechar{λ}{\ensuremath{\lambda}}
\newunicodechar{δ}{\ensuremath{\delta}}
\newunicodechar{Σ}{\ensuremath{\Sigma}}

% --- Colores institucionales ------------------------------------------
\definecolor{guinda}{HTML}{6C1D45}
\definecolor{azulescom}{HTML}{006699}
\definecolor{gristenue}{HTML}{F4F4F6}

% --- Títulos ----------------------------------------------------------
\titleformat{\section}{\normalfont\large\bfseries\color{guinda}}{\thesection.}{0.6em}{}
\titleformat{\subsection}{\normalfont\normalsize\bfseries}{\thesubsection.}{0.6em}{}
\titleformat{\subsubsection}{\normalfont\normalsize\bfseries\itshape}{\thesubsubsection.}{0.6em}{}
\titlespacing*{\section}{0pt}{1.2em}{0.4em}
\titlespacing*{\subsection}{0pt}{0.9em}{0.3em}
\titlespacing*{\subsubsection}{0pt}{0.7em}{0.2em}

% --- Encabezado y pie -------------------------------------------------
\pagestyle{fancy}
\fancyhf{}
\fancyhead[L]{\small\textcolor{guinda}{Teoría de la Computación}}
\fancyhead[R]{\small\PracticaCorta}
\fancyfoot[C]{\small\thepage}
\renewcommand{\headrulewidth}{0.4pt}

% --- Tablas y listas que genera pandoc --------------------------------
\providecommand{\tightlist}{\setlength{\itemsep}{0pt}\setlength{\parskip}{0pt}}
\providecommand{\passthrough}[1]{#1}
\AtBeginEnvironment{longtable}{\singlespacing\small}
\setlength{\LTpre}{0.6em}\setlength{\LTpost}{0.6em}
\renewcommand{\arraystretch}{1.25}

% --- Código (resaltado de pandoc, estilo tango) -----------------------
\usepackage{fancyvrb}
\usepackage{fvextra}
\newcommand{\VerbBar}{|}
\newcommand{\VERB}{\Verb[commandchars=\\\{\}]}
\DefineVerbatimEnvironment{Highlighting}{Verbatim}{commandchars=\\\{\},fontsize=\small,breaklines,breakanywhere}
\usepackage{framed}
\colorlet{shadecolor}{gristenue}
\newenvironment{Shaded}{\begin{singlespace}\begin{snugshade}}{\end{snugshade}\end{singlespace}}
\newcommand{\AlertTok}[1]{\textcolor[rgb]{0.94,0.16,0.16}{#1}}
\newcommand{\AnnotationTok}[1]{\textcolor[rgb]{0.56,0.35,0.01}{\textbf{\textit{#1}}}}
\newcommand{\AttributeTok}[1]{\textcolor[rgb]{0.13,0.29,0.53}{#1}}
\newcommand{\BaseNTok}[1]{\textcolor[rgb]{0.00,0.00,0.81}{#1}}
\newcommand{\BuiltInTok}[1]{#1}
\newcommand{\CharTok}[1]{\textcolor[rgb]{0.31,0.60,0.02}{#1}}
\newcommand{\CommentTok}[1]{\textcolor[rgb]{0.56,0.35,0.01}{\textit{#1}}}
\newcommand{\CommentVarTok}[1]{\textcolor[rgb]{0.56,0.35,0.01}{\textbf{\textit{#1}}}}
\newcommand{\ConstantTok}[1]{\textcolor[rgb]{0.56,0.35,0.01}{#1}}
\newcommand{\ControlFlowTok}[1]{\textcolor[rgb]{0.13,0.29,0.53}{\textbf{#1}}}
\newcommand{\DataTypeTok}[1]{\textcolor[rgb]{0.13,0.29,0.53}{#1}}
\newcommand{\DecValTok}[1]{\textcolor[rgb]{0.00,0.00,0.81}{#1}}
\newcommand{\DocumentationTok}[1]{\textcolor[rgb]{0.56,0.35,0.01}{\textbf{\textit{#1}}}}
\newcommand{\ErrorTok}[1]{\textcolor[rgb]{0.64,0.00,0.00}{\textbf{#1}}}
\newcommand{\ExtensionTok}[1]{#1}
\newcommand{\FloatTok}[1]{\textcolor[rgb]{0.00,0.00,0.81}{#1}}
\newcommand{\FunctionTok}[1]{\textcolor[rgb]{0.13,0.29,0.53}{\textbf{#1}}}
\newcommand{\ImportTok}[1]{#1}
\newcommand{\InformationTok}[1]{\textcolor[rgb]{0.56,0.35,0.01}{\textbf{\textit{#1}}}}
\newcommand{\KeywordTok}[1]{\textcolor[rgb]{0.13,0.29,0.53}{\textbf{#1}}}
\newcommand{\NormalTok}[1]{#1}
\newcommand{\OperatorTok}[1]{\textcolor[rgb]{0.81,0.36,0.00}{\textbf{#1}}}
\newcommand{\OtherTok}[1]{\textcolor[rgb]{0.56,0.35,0.01}{#1}}
\newcommand{\PreprocessorTok}[1]{\textcolor[rgb]{0.56,0.35,0.01}{\textit{#1}}}
\newcommand{\RegionMarkerTok}[1]{#1}
\newcommand{\SpecialCharTok}[1]{\textcolor[rgb]{0.81,0.36,0.00}{\textbf{#1}}}
\newcommand{\SpecialStringTok}[1]{\textcolor[rgb]{0.31,0.60,0.02}{#1}}
\newcommand{\StringTok}[1]{\textcolor[rgb]{0.31,0.60,0.02}{#1}}
\newcommand{\VariableTok}[1]{\textcolor[rgb]{0.00,0.00,0.00}{#1}}
\newcommand{\VerbatimStringTok}[1]{\textcolor[rgb]{0.31,0.60,0.02}{#1}}
\newcommand{\WarningTok}[1]{\textcolor[rgb]{0.56,0.35,0.01}{\textbf{\textit{#1}}}}
\RecustomVerbatimEnvironment{verbatim}{Verbatim}{fontsize=\small,breaklines,breakanywhere}

% --- Evidencias -------------------------------------------------------
% \evidencia[ancho]{ruta}{leyenda}: inserta la captura si existe; si aún
% no se ha agregado, dibuja un recuadro que indica la ruta esperada.
\newcommand{\evidencia}[3][0.95]{%
  \begin{figure}[htbp]\centering
  \IfFileExists{#2}{%
    \fbox{\includegraphics[width=#1\linewidth,height=0.42\textheight,keepaspectratio]{#2}}%
  }{%
    \fcolorbox{guinda}{guinda!6}{\parbox[c][3.2cm][c]{0.8\linewidth}{\centering\singlespacing
      \textbf{\textcolor{guinda}{Captura pendiente}}\\[2pt]
      \footnotesize\ttfamily\detokenize{#2}}}%
  }
  \caption{#3}
  \end{figure}}
% \evidenciaaqui: igual que \evidencia, pero sin flotar (se queda junto al texto que la explica).
\newcommand{\evidenciaaqui}[3][0.95]{%
  \par\begin{center}
  \IfFileExists{#2}{%
    \fbox{\includegraphics[width=#1\linewidth,height=0.4\textheight,keepaspectratio]{#2}}%
  }{%
    \fcolorbox{guinda}{guinda!6}{\parbox[c][3.2cm][c]{0.8\linewidth}{\centering\singlespacing
      \textbf{\textcolor{guinda}{Captura pendiente}}\\[2pt]
      \footnotesize\ttfamily\detokenize{#2}}}%
  }
  \captionof{figure}{#3}
  \end{center}}
